\chapter{Resultados e Discussões}
\thispagestyle{headings}

    \par As implicações dessas alterações são profundas, tanto em termos de implementações novas quanto de manutenibilidade. No que tange à implementação, a nova abordagem exige uma estrutura de código mais elaborada. A criação de uma única funcionalidade, que antes se resumia a um método em um Controller, agora envolve a definição de Entidades, Casos de Uso com seus respectivos DTOs e Boundaries, e as implementações dos adaptadores, como o Repository e o Presenter.
    
    \par É notável que embora isso represente um aumento na quantidade de arquivos e na complexidade inicial, a implementação passa a ser guiada por contratos (interfaces), reforçando o Princípio da Inversão de Dependência e resultando em um código mais explícito e coeso. É na manutenibilidade, no entanto, que os maiores benefícios serão percebidos, sendo este um dos objetivos centrais da Arquitetura Limpa. Ao isolar o núcleo da aplicação (Casos de Uso e Entidades) dos detalhes de infraestrutura como o \textit{framework} e o banco de dados, a arquitetura facilita a manutenção futura. 
    
    \par Na prática, a substituição do Eloquent ORM por outra ferramenta de persistência, por exemplo, exigiria apenas a criação de uma nova classe que implemente a ContactRepositoryInterface, sem qualquer alteração na lógica de negócio contida no Interactor. Essa separação estrita de interesses não só protege as regras de negócio de mudanças em tecnologias externas, mas também pode tornar o sistema mais flexível e testável.
    
    \par Contudo, é fundamental aprofundar a discussão sobre os \textit{trade-offs} (compensações) envolvidos nessa adoção, que vão além da implementação técnica e afetam diretamente a dinâmica da equipe de desenvolvimento, o ciclo de vida do projeto e a alocação de recursos.
    
    \par Um desses desafios é a curva de aprendizado e a integração da equipe. Na adoção da Arquitetura Limpa em um ecossistema como o Laravel, a dificuldade imposta ao time é evidente. A estrutura nativa do \textit{framework} é amplamente conhecida e documentada, e sua filosofia de ``convenção sobre configuração'' permite que desenvolvedores, mesmo com menos experiência, se tornem produtivos rapidamente e saibam onde encontrar cada componente do sistema. 
    
    \par Contudo, essa facilidade inicial não garante que o código esteja bem organizado: pois pode ser possível encontrar regras de negócio profundamente acopladas aos \textit{Controllers}, validações misturadas com lógica de apresentação, e dependências diretas do \textit{Eloquent} espalhadas por toda a aplicação. Principalmente em projetos onde não houve uma preocupação com arquitetura desde o inicio.
    
    \par Ao adotar a Arquitetura Limpa, o projeto se desvia significativamente dessa convenção em favor de uma organização mais rígida e explícita. Profissionais com pouco ou nenhum conhecimento prévio em \textit{design} de \textit{software} orientado a camadas, SOLID e padrões de projeto (como \textit{Repository}, DTOs e \textit{Data Mappers}) podem encontrar uma dificuldade inicial considerável.

    \par Isso implica que a integração de novos membros à equipe (processo de \textit{onboarding}) exigirá um período de treinamento e mentoria mais longo. A equipe precisa não apenas entender o domínio de negócio, mas também a filosofia arquitetural do projeto, que agora possui suas próprias convenções, independentes do \textit{framework}.
    
    \par Outro ponto de análise é o custo de desenvolvimento de novas funcionalidades. Conforme apontado, a complexidade inicial e o aumento no número de arquivos refletem-se diretamente no tempo de desenvolvimento. A implementação de uma nova tarefa ou melhoria, por menor que seja, exigirá um esforço inicial maior. 
    
    \par A criação de um novo \textit{endpoint} de CRUD, por exemplo, que no MVC padrão do Laravel poderia ser resumida a um ou dois métodos no \textit{Controller} e ajustes no \textit{Model}, agora demanda a criação de uma cadeia de artefatos: a definição (ou reutilização) da Entidade, a criação do Caso de Uso (\textit{Interactor}), a definição das \textit{Boundaries} (entrada e saída), a criação dos DTOs (\textit{Request} e \textit{Response}), a implementação do \textit{Controller} e do \textit{Presenter} (como Adaptadores), a atualização da interface do Repositório (se necessário) e, por fim, a implementação no Repositório concreto.
    
    \par Esse ``\textit{boilerplate}'' (código repetitivo) de infraestrutura arquitetural, embora garanta o desacoplamento, inevitavelmente torna o ciclo de desenvolvimento de novas \textit{features} mais lento. Este é um custo que deve ser conscientemente aceito pela equipe e pelo negócio em troca dos benefícios de longo prazo.
       
    \par Deve-se considerar, também, a aplicação em projetos legados e de alta complexidade. O presente estudo de caso partiu de uma API REST base relativamente simples. A transposição dessa metodologia para um sistema de \textit{software} em produção, especialmente um projeto legado ou de alta complexidade, é um desafio de outra magnitude.
    
    \par Em sistemas legados, o acoplamento profundo das regras de negócio com o \textit{framework} (especialmente com o Eloquent) é comum. A tentativa de refatorar um sistema nessas condições pode ser paralisante e introduzir um alto risco de geração de novos bugs. Nesses cenários, a implementação da Arquitetura Limpa provavelmente exigiria uma equipe dedicada e uma estratégia de migração gradual, onde novas funcionalidades são construídas seguindo a nova arquitetura e funcionalidades antigas são refatoradas incrementalmente. O custo e o tempo para essa transição seriam significativamente maiores, podendo justificar a alocação de uma equipe específica apenas para essa modernização arquitetural.
    
    \par Em contrapartida direta ao aumento do custo de implementação, a arquitetura oferece vantagens claras na manutenibilidade, testabilidade e na resolução de bugs. A principal vantagem reside na testabilidade aprimorada, um dos pilares da arquitetura.
    
    \par O isolamento rigoroso do \textit{Interactor} (camada de Casos de Uso) permite a criação de testes unitários robustos que validam a lógica de negócio pura, sem a necessidade de simular o \textit{framework}, requisições HTTP ou conexões com o banco de dados. As dependências externas (como o repositório) são injetadas via interfaces, permitindo o uso de dublês de teste de forma mais fácil.
    
    \par Essa cobertura de testes na camada mais crítica do sistema facilita a rápida identificação e localização precisa de falhas. Quando um \textit{bug} é reportado, pode ser mais fácil diagnosticá-lo: a falha está na tradução de dados feita pelo \textit{Controller}? Está na formatação do \textit{Presenter}? Ou está na regra de negócio dentro do \textit{Interactor}? Como o \textit{Interactor} é puramente testável, ele pode ser validado ou descartado como fonte do problema rapidamente, tornando o processo de \textit{debugging} mais eficiente e rápido.

    \par Em suma, para que uma equipe de desenvolvimento ou uma empresa decida adotar a Arquitetura Limpa, é fundamental avaliar o \textit{trade-off} entre o custo inicial de complexidade e o ganho de longevidade. Essa análise de viabilidade deve ser ancorada em três fatores críticos: a complexidade do projeto, o tamanho e a experiência da equipe, e a necessidade de mudança.

    \par O primeiro fator é a complexidade do projeto. A arquitetura é idealmente indicada para sistemas de alta complexidade, a qual se caracteriza por ter alta volatilidade nas regras de negócio e uma expectativa de longevidade a longo prazo (tipicamente, cinco anos ou mais). Para projetos de escopo simples, que demandam agilidade imediata, o custo de \textit{boilerplate} não se justifica. Nesses casos, o benefício da manutenibilidade é ofuscado pelo alto custo de desenvolvimento inicial, pois o tempo gasto na criação das múltiplas camadas de adaptação e interfaces consome recursos que seriam melhor utilizados na entrega rápida de \textit{features} através de um modelo mais simples, como o MVC nativo.
    
    \par O segundo fator é o tamanho e a experiência da equipe. Equipes grandes, ou aquelas que enfrentam alta rotatividade de membros, apresentam um cenário particularmente favorável à Arquitetura Limpa. Embora o processo de \textit{onboarding} demande um tempo inicial maior para compreensão dos conceitos arquiteturais, a clareza estrutural compensa essa complexidade ao oferecer um mapa mental consistente do sistema.
    
    \par Um exemplo disso pode ser a saída de um membro experiente, que não ``leva'' consigo a organização e as regras de funcionamento de uma parte do sistema desenvolvida por ele, tornando a substituição de membros mais segura, ainda que exija treinamento técnico prévio.
    
    \par Diferentemente de arquiteturas mais flexíveis (onde convenções podem variar entre módulos), as fronteiras rígidas e explícitas da Arquitetura Limpa reduzem a ambiguidade: um novo desenvolvedor sabe que regras de negócio estão nos Interactors, adaptações de dados nos Presenters, e assim por diante. Essa previsibilidade, uma vez internalizada, facilita a localização de código, a compreensão de fluxos e a realização de mudanças com maior segurança, mesmo sem profundo conhecimento do domínio de negócio.

    \par Em contrapartida, equipes pequenas podem não se beneficiar tanto desse ganho de clareza estrutural. Como a comunicação entre poucos desenvolvedores tende a ser mais direta e o conhecimento do sistema pode ser compartilhado organicamente através de conversas e revisões de código, o benefício da estrutura rígida e auto-explicativa é menos significativo. 
    
    \par Nesses casos, o alto custo de desenvolvimento inicial pode não compensar, especialmente quando a equipe já possui canais eficientes de comunicação e compartilhamento de conhecimento. A perda de agilidade a curto prazo torna-se um fator mais crítico, pois o \textit{overhead} arquitetural consome recursos que, em uma equipe reduzida, fazem mais falta na entrega de funcionalidades.

    \par Em qualquer cenário de tamanho de equipe, o modelo exige que a mesma possua experiência consolidada em princípios de design (como SOLID e Padrões de Projeto), pois o modelo remove a “mágica” do framework em favor de um código mais explícito, o que pode aumentar a curva de aprendizado para desenvolvedores menos experientes.
    
    \par Por fim, a necessidade de mudança e o tempo disponível são cruciais. A necessidade de mudança é o principal \textit{driver} de implementação, pois representa a urgência em proteger o software de futuras alterações tecnológicas e garantir sua longevidade. Contudo, o tempo disponível atua como o fator limitante. Se a necessidade de longevidade for baixa, ou se não houver tempo e recursos adequados para absorver o custo de um ciclo de desenvolvimento inicial mais lento (o \textit{boilerplate}), a adoção da Arquitetura Limpa é desaconselhada.